\hypertarget{appendix-8836553d56fc}{}
\section*{Implicit differentiation and parameter sensitivities}
\section{Implicit differentiation and parameter sensitivities}\label{implicit-differentiation-and-parameter-sensitivities}

\textbf{Curriculum depth:} Developed · \textbf{First seen:} Phase 9 · \textbf{Latest development:} Phase 9

\subsection{Definition and mechanism}\label{definition-and-mechanism}

Implicit differentiation obtains how a solution changes when it is defined by an equation rather than an explicit formula. If a state satisfies F(y,θ)=0 and the Jacobian with respect to y is invertible, differentiating gives dy/dθ=−(∂F/∂y)⁻¹(∂F/∂θ). This can differentiate equilibrium states, constrained optima, and roots without unrolling every solver iteration. For a converged optimizer, differentiating stationarity and constraint conditions yields sensitivity equations, but active-set changes, nonunique solutions, singular Jacobians, and imperfect convergence complicate the result. Continuous sensitivity equations propagate derivatives alongside an ODE; adjoint methods can reduce cost when there are many parameters and few scalar outputs. These are modeling tools with assumptions, not a guarantee that every algorithmic branch is differentiable.

\subsection{Prerequisites}\label{prerequisites}

\begin{itemize}
\tightlist
\item
  \href{}{Automatic differentiation and backpropagation (03-Visual-Computing-I, appendix-40df0855b4a7)}
\item
  \href{}{Constrained optimization and inverse kinematics (04-Visual-Computing-II, appendix-29f4bea4c2bb)}
\item
  \href{}{ODE integration and stability (04-Visual-Computing-II, appendix-18eacf61a04b)}
\end{itemize}

\subsection{Used by}\label{used-by}

\begin{itemize}
\tightlist
\item
  This node is a current endpoint in the supplied Phase 9--10 path.
\end{itemize}

\subsection{Full teaching explanation}\label{full-teaching-explanation}

\begin{itemize}
\tightlist
\item
  \href{Continuous\%20Notes/Computing\%20Foundations\%20-\%20From\%20Bits\%20to\%20C++\%20and\%20Systems\%20APIs\#Numerical\%20differentiation\%20and\%20automatic\%20differentiation}{Phase 9 teaching section}
\item
  \href{}{Visual Recall/09 - Calculus and Optimization - Visual Recall.md (01-Foundations-I, appendix-adf48f5dc5c1)} · visual recall
\end{itemize}

\subsection{Source coverage}\label{source-coverage}

Supplied lesson parts: 855 (T03 · M052), 899 (T02 · M054), 900 (T02 · M054).
